% !TEX root = ../main.tex
\section{Spectrum and Sensitivity}
\label{sec:spectrum}
\label{sec:sensitivity}

% Combines the pedagogical ``what features appear in the spectrum'' discussion with
% the measurement-strategy / sensitivity arguments. Target ~7--8 pp total.

\subsection{Probability versus $L/E$}
% Global $P(L/E)$ --- fast atmospheric wiggle + slow solar envelope --- one
% paragraph and one figure.

\subsection{The NOvA Oscillation Window}
% Zoom to NOvA ($L=810$\,km, $E\sim 1$--$3$\,GeV): a narrow $L/E$ slice spanning
% one atmospheric oscillation cycle, not multiple wiggles. Physics motivation for
% this baseline and energy: sits at the oscillation maximum for $\Delta m^{2}_{32}$
% and gives a matter effect large enough to help break the ordering--$\delta_{\rm CP}$
% degeneracy. Numbers only; no hardware.

\subsection{The Disappearance Spectrum}
% FD disappearance: single broad dip --- position fixes $\Delta m^{2}_{32}$,
% depth fixes $\sin^{2}(2\theta_{23})$. Why the dip region ($L/E\sim L_{\rm osc}/2$)
% gives the sharpest $\Delta m^{2}_{32}$ handle.

\subsection{The Appearance Spectrum}
% FD appearance: single broad peak --- rate carries $\theta_{13}\times\delta_{\rm CP}$
% interference and (via matter) mass ordering; shape carries less information.
% $\nu$ vs $\bar\nu$ running.
% Appearance vs disappearance probe orthogonal parameter combinations.
% $\delta_{\rm CP}$ appears only in appearance (interference term);
% $\nu_e$/$\bar\nu_e$ rate difference is the primary CP handle.
% Small derivation of matter-modified $P(\nu_\mu\to\nu_e)$ (Nunokawa--Parke--Valle
% expansion; not a full three-eigenvalue diagonalisation). Consequence:
% $\nu$/$\bar\nu$ asymmetry combines matter + $\delta_{\rm CP}$ degenerately;
% a long baseline separates them because matter scales differently with $E$.
